\documentclass[12pt,a4paper]{article}

\usepackage[T1]{fontenc}
\usepackage[utf8]{inputenc}
\usepackage[turkish]{babel}
\usepackage{geometry}
\usepackage{booktabs}
\usepackage{tabularx}
\usepackage{xcolor}
\usepackage{hyperref}
\usepackage{listings}
\usepackage{enumitem}
\usepackage{parskip}

\geometry{
  a4paper,
  top=2.5cm,
  bottom=2.5cm,
  left=3cm,
  right=2.5cm
}

\definecolor{codeblue}{RGB}{32,78,145}
\definecolor{codegray}{RGB}{245,246,248}
\definecolor{codegreen}{RGB}{35,120,80}
\definecolor{codered}{RGB}{170,45,45}

\lstset{
  backgroundcolor=\color{codegray},
  basicstyle=\ttfamily\small,
  keywordstyle=\color{codeblue}\bfseries,
  commentstyle=\color{codegreen}\itshape,
  stringstyle=\color{codered},
  numbers=left,
  numberstyle=\tiny,
  frame=single,
  breaklines=true,
  showstringspaces=false,
  tabsize=2
}

\hypersetup{
  colorlinks=true,
  linkcolor=codeblue,
  urlcolor=codeblue,
  pdftitle={CodeAlchemist VS Code Extension Hafta 14},
  pdfauthor={CodeAlchemist Ekibi}
}

\title{\textbf{CodeAlchemist VS Code Extension --- Hafta 14} \\
\large Agent Baglam Koruma, Komut Ciktisi ve Saglik Durumu Stabilizasyonu}
\author{CodeAlchemist Ekibi}
\date{18 Mayis 2026}

\begin{document}

\maketitle
\tableofcontents
\newpage

\section{Yonetici Ozeti}

Bu hafta VS Code extension tarafinda agent deneyimini mevcut modern kod asistanlarina yaklastiran uc kritik alan iyilestirildi:

\begin{itemize}[leftmargin=2em]
  \item Sohbet icinde dosya hedefi ve calisma baglami korunur hale getirildi.
  \item Agent tarafindan onerilen terminal komutlari yakalanabilir sekilde calistirilip stdout/stderr ciktisi chat kartina yansitildi.
  \item Backend saglik durumu ve auth handshake akisi tekil ping hatalarinda UI'yi yanlis offline durumuna dusurmemek icin sertlestirildi.
\end{itemize}

\section{Kod Tabani Kapsami}

\begin{tabularx}{\textwidth}{@{}lX@{}}
\toprule
\textbf{Dosya} & \textbf{Degisiklik} \\
\midrule
\texttt{src/chatProvider.ts} & Sohbet bazli working context, komut calistirma, health/auth yeniden yayinlama ve aksiyon karti akisi guncellendi. \\
\texttt{src/actionHandler.ts} & Tek cevapta dosya degisikligi ve komut aksiyonunun birlikte yakalanmasi icin komut parser'i eklendi. \\
\texttt{src/chatView.ts} & Komut karti ciktisi, webview mesaj dinleyicisi ve uzun terminal log gorunumu guncellendi. \\
\texttt{src/healthMonitor.ts} & Tekil health ping hatalarina tolerans eklendi; basarili auth durumunda online state korunur hale getirildi. \\
\texttt{src/types.ts} & \texttt{client\_context.working\_context} ve auth/limit alanlari tip sistemine eklendi. \\
\bottomrule
\end{tabularx}

\section{Baglam Koruma}

Onceki davranista extension her istekte yalnizca aktif editor ve workspace snapshot'ini gonderiyordu. Kullanici ikinci mesajda ``bunu duzelt'', ``ayni dosyaya ekle'' veya ``devam et'' dediginde model dosya hedefini kaybedebiliyor ve tekrar dosya adi sorabiliyordu.

Yeni akista \texttt{chatProvider.ts} icinde sohbet anahtarina bagli bir working context tutulur:

\begin{itemize}[leftmargin=2em]
  \item aktif dosya,
  \item son hedef dosya,
  \item acikca adi gecen dosyalar,
  \item agent tarafindan son degistirilen dosyalar.
\end{itemize}

Bu bilgiler her istekte \texttt{client\_context.working\_context} alanina ve prompt icindeki continuity instruction bolumune eklenir. Boylece kullanici ayni sohbet icinde takip mesaji yazdiginda hedef dosya korunur.

\section{Terminal Komutu Ciktisi}

VS Code Terminal API terminale komut gondermeye izin verir; fakat terminal stdout/stderr ciktisini extension tarafina geri okumak icin guvenilir bir API sunmaz. Bu nedenle komutlar artik extension host icinde \texttt{child\_process.spawn} ile calistirilmaktadir.

\begin{lstlisting}[language=TypeScript, caption={Komut ciktisinin webview'e aktarilmasi}]
child.stdout?.on('data', (data: Buffer) => {
  this._postCommandOutput(actionId, data.toString('utf8'), 'stdout');
});

child.stderr?.on('data', (data: Buffer) => {
  this._postCommandOutput(actionId, data.toString('utf8'), 'stderr');
});
\end{lstlisting}

Windows ortaminda komutlar PowerShell uzerinden, Unix tabanli ortamlarda ise \texttt{/bin/bash -lc} uzerinden calistirilir. Calisma dizini workspace disina tasirilmadan once dogrulanir.

\section{Coklu Aksiyon Duzeltmesi}

Agent cevaplari bazen ayni turda hem dosya degisikligi hem terminal komutu icerebilir. Eski parser ilk dosya degisikligini bulunca durdugu icin komut karti olusmuyordu. \texttt{parseCommandActions(...)} yardimcisi ile trace, structured response ve embedded JSON icindeki \texttt{run\_command} aksiyonlari ayrica taranmaya baslandi.

Beklenen UI davranisi:

\begin{enumerate}[leftmargin=2em]
  \item Dosya degisikligi icin \textbf{File Change} karti gosterilir.
  \item Komut icin ayrica \textbf{Terminal Komutu} karti gosterilir.
  \item Kullanici onay verdiginde komut ciktisi ayni kart icinde canli olarak akar.
\end{enumerate}

\section{Saglik Durumu ve Auth Handshake}

Gozlenen problemde backend \texttt{/v1/status} istegine basarili cevap verirken, tekil \texttt{/health} ping hatasi UI'yi \texttt{offline} veya \texttt{connecting} durumuna cekebiliyordu. Bu akista iki sertlestirme yapildi:

\begin{itemize}[leftmargin=2em]
  \item \texttt{HealthMonitor} ardisik iki health hatasi olmadan mevcut online durumunu offline'a dusurmez.
  \item Basarili auth status kontrolu \texttt{markOnline(...)} ile health state'i tekrar online kabul eder.
\end{itemize}

Bu sayede backend gercekte erisilebilirken UI'nin yalanci offline durumuna dusmesi engellenir.

\section{Dogrulama}

Degisikliklerden sonra extension paketi TypeScript derlemesinden gecirildi:

\begin{lstlisting}[language=bash, caption={Derleme dogrulamasi}]
npm run build
\end{lstlisting}

Backend saglik kontrolu de yerelde dogrulandi:

\begin{lstlisting}[language=bash, caption={Backend health kontrolu}]
GET http://localhost:5001/health
{"status":"ok","version":"1.0.0"}
\end{lstlisting}

\section{Test Senaryosu}

Extension Development Host yeniden yuklendikten sonra asagidaki prompt ile hem dosya baglami hem komut ciktisi test edilebilir:

\begin{quote}
Bu aktif dosyanin sonuna "Komut testi" yaz. Sonra \texttt{Get-Date} komutunu calistir ve ciktisini chat icinde goster.
\end{quote}

Ikinci mesaj olarak:

\begin{quote}
Ayni dosyanin sonuna bir satir daha ekle.
\end{quote}

Bu testte agent'in tekrar dosya adi sormadan onceki hedef dosyayi kullanmasi beklenir.

\end{document}
